% ============================ 第一节 释义 ====================================
% 基本术语按主体分类（\sygroup 通栏分类行，同 2026 年申报稿）+ 专业术语；结尾放全书通用的四舍五入尾差说明（申报文件惯例）。
\gssection{释义}

本招股说明书中，除非文义另有所指，下列词语具有如下含义：

\subsection{基本术语}

\begin{shiyi}
\sygroup{发行人}
\sy{发行人、公司、本公司、示例科技}{示例科技股份有限公司}
\sy{示例有限}{示例科技有限公司，发行人前身}
\sygroup{发行人子公司}
\sy{示例智算}{示例智算科技有限公司，发行人全资子公司}
\sy{示例数据}{示例数据科技有限公司，发行人全资子公司}
\sygroup{发行人股东及员工持股平台}
\sy{控股股东、示例控股}{示例控股有限公司}
\sy{实际控制人}{林××}
\sy{示例聚力}{示例市示例聚力投资合伙企业（有限合伙），发行人员工持股平台}
\sy{示例创投}{示例创业投资基金合伙企业（有限合伙），发行人股东}
\sy{示例产投}{示例产业投资基金合伙企业（有限合伙），发行人股东}
\sygroup{本次发行及中介机构}
\sy{本次发行}{发行人首次公开发行12,000.00万股人民币普通股（A股）并在上海证券交易所科创板上市的行为}
\sy{保荐人、主承销商、示例证券}{示例证券股份有限公司}
\sy{申报会计师、示例会计师}{示例会计师事务所（特殊普通合伙）}
\sy{发行人律师}{示例市示例律师事务所}
\sygroup{政府机构和部门}
\sy{中国证监会}{中国证券监督管理委员会}
\sy{上交所、交易所}{上海证券交易所}
\sygroup{法律法规、公司制度及治理机构}
\sy{《公司法》}{《中华人民共和国公司法》}
\sy{《证券法》}{《中华人民共和国证券法》}
\sy{《公司章程》}{《示例科技股份有限公司章程》}
\sy{《公司章程（草案）》}{《示例科技股份有限公司章程（草案）》，自公司股票在上交所科创板上市之日起生效}
\sy{股东会}{示例科技股份有限公司股东会}
\sy{董事会}{示例科技股份有限公司董事会}
\sy{审计委员会}{示例科技股份有限公司董事会审计委员会。根据《公司法》及《公司章程》的规定，公司不设监事会，由审计委员会行使《公司法》规定的监事会职权}
\sygroup{其他}
\sy{报告期}{2023年度、2024年度、2025年度}
\sy{报告期各期末}{2023年12月31日、2024年12月31日、2025年12月31日}
\sy{高新技术企业}{依据《高新技术企业认定管理办法》认定的高新技术企业}
\sy{A股}{在中国境内证券交易所上市、以人民币标明面值、以人民币认购和交易的人民币普通股}
\sy{元、万元、亿元}{人民币元、人民币万元、人民币亿元}
\end{shiyi}

\subsection{专业术语}

\begin{shiyi}
\sy{人工智能、AI}{Artificial Intelligence，利用计算机模拟、延伸和扩展人类智能的理论、方法、技术及应用系统}
\sy{大模型、基础模型}{基于海量数据预训练、参数规模庞大、具备通用理解与生成能力，并可通过微调等方式适配下游任务的深度学习模型}
\sy{Transformer}{基于自注意力机制的深度神经网络架构，是当前大模型的主流基础架构}
\sy{MoE、混合专家模型}{Mixture of Experts，将模型划分为多个专家子网络并按输入动态激活部分专家的模型架构，可在控制推理成本的同时扩大模型参数规模}
\sy{多模态}{融合文本、图像、音频、视频等多种信息形态的模型能力，包括多模态理解与多模态生成}
\sy{预训练}{在大规模无标注语料上以自监督方式训练模型基础能力的过程}
\sy{微调、SFT}{Supervised Fine-Tuning，监督微调，使用标注数据对预训练模型进行下游任务适配训练的过程}
\sy{RLHF}{Reinforcement Learning from Human Feedback，基于人类反馈的强化学习，使模型输出与人类偏好和价值观对齐的训练方法}
\sy{推理}{已训练完成的模型接收输入并生成输出的计算过程，即模型对外提供服务的运行阶段}
\sy{Token}{大模型处理文本的基本单元，开放平台服务通常按输入、输出Token数量计量计费}
\sy{MaaS}{Model as a Service，模型即服务，通过开放平台以API等方式向客户提供模型能力并按用量计费的服务模式}
\sy{API}{Application Programming Interface，应用程序编程接口}
\sy{智能体、Agent}{以大模型为核心，具备任务规划、工具调用、记忆与自主执行能力的智能应用形态}
\sy{RAG}{Retrieval-Augmented Generation，检索增强生成，将外部知识检索与模型生成相结合的技术}
\sy{AIGC}{Artificial Intelligence Generated Content，人工智能生成内容}
\sy{算力}{支撑模型训练与推理的智能计算资源，通常以GPU等智能计算芯片集群形式提供}
\sy{私有化部署}{将模型及配套系统部署于客户本地或专属云环境，由客户在其管理域内使用的交付方式}
\end{shiyi}

本招股说明书中若出现总数与各分项数值之和尾数不符的情况，均为四舍五入所致。
